\documentclass[12pt]{article}
\usepackage{amsmath,amssymb,amsthm}
\usepackage{geometry}
\usepackage{hyperref}
\usepackage{natbib}
\usepackage{enumerate}
\geometry{margin=1in}

\newtheorem{definition}{Definition}
\newtheorem{proposition}{Proposition}
\newtheorem{remark}{Remark}

\title{What Survives Compaction?\\
On the Missing Principle Behind Context Compression\\
in Reasoning Models}

\author{Minamo Minamoto\\
\small Independent\\
\small ORCID: 0009-0002-1201-5704\\
\small \href{https://github.com/minamominamoto/topology-of-grounding}{topology-of-grounding}
}

\date{April 2026}

\begin{document}
\maketitle

\begin{abstract}
Recent work on context compression for reasoning models---most notably
Memento \citep{kontonis2026memento}---has demonstrated that large language
models can learn to segment their chain-of-thought, compress each segment
into a dense summary, and reason forward from summaries alone, achieving
2--3$\times$ KV cache reduction with minimal accuracy loss on in-distribution
benchmarks. We argue that this technical achievement raises a question that
engineering metrics alone cannot answer: \emph{what should survive
compaction, and why?} Current compression approaches optimize for accuracy
preservation on downstream benchmarks, treating compaction as a performance
engineering problem. We propose that this framing is structurally incomplete.
Drawing on the interpretive stack model of cognition
\citep{minamoto2026overwrite,minamoto2026accum}, we argue that compaction is
fundamentally a frame-selection problem: the compression operation implicitly
decides which interpretive frames are foregrounded and which are archived.
Without a principled account of this decision, compaction risks
systematically suppressing minority frames that are essential for error
correction under distributional shift, while amplifying the dominant frame
regardless of its appropriateness to the current reasoning context.
We propose three principles for SSC-aware compression and discuss their
implications for the design of Memento-style, LightThinker++, and
neural-computer architectures.
\end{abstract}

\section{Introduction}

Context compression for reasoning models addresses a genuine engineering
challenge: chain-of-thought reasoning generates long intermediate sequences
that consume KV cache, increase inference latency, and scale poorly with
problem depth. Memento \citep{kontonis2026memento} addresses this by training
models to segment their reasoning and replace each completed segment with a
compact summary (``memento''), from which subsequent reasoning proceeds.
The approach is technically elegant and achieves impressive compression
ratios with minimal benchmark degradation.

We do not dispute the technical contribution. We argue, however, that the
framing of compaction as a performance engineering problem---optimize
compression ratio subject to accuracy constraints---misses a structural
question that has consequences for the robustness of compressed systems
under distribution shift.

The question is: what should survive compaction?

Current systems implicitly answer: whatever maximizes accuracy on held-out
test sets from the same distribution as training. We argue this answer is
insufficient because it conflates task-accurate compression with
\emph{cognitively safe} compression. The two can diverge precisely in the
situations where compression matters most: long, complex reasoning chains
that approach or exceed the model's in-context capacity, and that may
involve reasoning under uncertainty where multiple interpretive frames are
relevant.


\begin{remark}[Terminological note on ``overwrite'']
See \cite{minamoto2026overwrite} for the precise definition.
In this paper, \emph{overwrite} denotes a frame-dominant transition
(superposition $+$ selection-dominance), not deletion of prior frames.
\end{remark}

\section{The Interpretive Stack Model: Brief Summary}

We summarize the relevant framework from \citet{minamoto2026overwrite,minamoto2026accum};
readers familiar with these papers may skip to Section~3.

\begin{definition}[Interpretive Frame]
An interpretive frame $F$ is a function from inputs to outputs induced
by a specific configuration of a language model's weight space---
operationally, a cluster of circuits whose causal contribution to output
is coherent across a class of inputs.
\end{definition}

\begin{definition}[Interpretive Stack]
An interpretive stack $L = (F_1, F_2, \ldots, F_n)$ is an ordered
sequence of interpretive frames implicitly encoded in the model's weight
space. $F_1$ is the most recently installed or highest-priority frame.
Earlier frames are residual structures from pretraining or prior
fine-tuning stages.
\end{definition}

\begin{definition}[Selection Stability Capacity (SSC)]
A model exhibits Selection Stability Capacity with respect to a target
frame $F^*$ if the selection rule $S$ assigns $F^*$ across the full
distribution of intended input contexts. SSC failure occurs when $S$
activates a frame $F_k \neq F^*$ for some input in the intended distribution.
\end{definition}

The central claim of \citet{minamoto2026accum} is that prior interpretive
frames are not overwritten by fine-tuning but archived: they remain
accessible and may be reactivated under appropriate conditions. SSC
measures the model's ability to maintain the correct frame despite the
presence of archived alternative frames.

\section{Compaction as a Frame-Selection Problem}

When a compressor summarizes a reasoning block, it necessarily selects
what to preserve. In the absence of an explicit principle, this selection
is determined by the compressor's training signal---accuracy on
in-distribution follow-up tasks.

We argue this training signal systematically biases compression toward
preserving the dominant frame's conclusions while discarding information
about alternative frames that were considered and provisionally rejected.
The result is a memento that encodes the correct answer for in-distribution
inputs but loses the structural information needed to recover when the
subsequent reasoning context shifts to a regime where the suppressed
alternative frame would be more appropriate.

This is not a hypothetical concern. It is the structural consequence of
optimizing compression on i.i.d.\ held-out benchmarks: if the training
distribution is well-covered by the dominant frame, the compressor will
never receive a gradient signal to preserve alternative-frame indices,
because those indices are never needed on training inputs.

\section{Three Principles for SSC-Aware Compression}

We propose three principles that a compaction approach should satisfy to
be cognitively safe.

\textbf{Principle 1: Conclusion preservation (existing).}
The compressed representation should encode the conclusion reached at
the end of the compressed block. This is what Memento already does, and
it correctly captures the minimal information needed for in-distribution
continuation.

\textbf{Principle 2: Alternative-frame index preservation (missing).}
The compressed representation should encode, for each alternative frame
that was considered and rejected during the compressed block:
\begin{enumerate}[(i)]
\item The identity of the alternative frame (or a compact representation
  of its acquisition context);
\item The condition under which it was rejected (the evidence that
  favored the dominant frame);
\item The trigger conditions under which it should be re-evaluated
  (input features that would shift the balance of evidence).
\end{enumerate}

This index is typically much smaller than the full frame content.
Preserving it does not require storing the alternative frame's full
reasoning trace---only the information needed to decide whether to
re-activate it.

\textbf{Principle 3: SSC-aware compression evaluation (missing).}
Compression quality should be evaluated not only on task accuracy under
the training distribution but on Selection Stability Capacity: the
compressed state's ability to maintain the correct frame under
distributional shift. A memento that achieves high accuracy on
in-distribution follow-up blocks but low SSC is brittle.

Operationally, SSC-aware evaluation would require:
\begin{enumerate}[(i)]
\item Deliberately perturbing follow-up reasoning contexts to resemble
  the acquisition context of the suppressed alternative frame.
\item Measuring whether the compressed model correctly maintains the
  dominant frame or erroneously reverts to the suppressed frame.
\item Using this measurement as an additional compression quality signal
  during compressor training.
\end{enumerate}

\section{Relation to LightThinker++, Neural Computer, and Thinking Mid-Training}

LightThinker++ \citep{lightthinkerplusplus2026} introduces three explicit
memory operations: Commit, Expand, and Fold. The Expand operation implicitly
acknowledges that archived content may need retrieval. However,
LightThinker++ does not address what should be archived beyond
task-relevant intermediate results. Under the interpretive stack model,
the Expand operation is insufficient if the archived content did not
include alternative-frame indices: expanding a memento that encodes only
the dominant frame's conclusion cannot recover the suppressed alternative.

The Neural Computer paradigm \citep{neuralcomputer2026} collapses
computation, memory, and I/O into a single latent runtime state.
This is architecturally elegant but does not resolve the frame-selection
problem---it relocates it to the question of what the latent state encodes.
The same compaction principle applies.

Meta FAIR's thinking mid-training approach \citep{lanchantin2026midtraining}
inserts an intermediate SFT+RL phase between pretraining and post-training,
achieving a 3.2$\times$ improvement in reasoning performance. The
interpretive stack model offers a structural explanation for this result:
mid-training installs a richer stack of interpretive frames before
post-training begins, increasing SSC at the time of post-training.
When post-training operates on a model with higher SSC, the resulting
fine-tuned frame is more stable under distributional shift.
Conversely, post-training without mid-training operates on a shallow
stack, producing a dominant frame with low SSC---the frame activates
correctly on i.i.d.\ inputs but fails under the modest perturbations
that mid-training examples provide.

This explanation connects the compaction problem to the training pipeline:
a system that enters compaction with high SSC (due to mid-training) will
produce mementos that are more structurally complete, because the model's
reasoning will more regularly surface alternative-frame considerations
rather than immediately collapsing to the dominant-frame conclusion.

\section{Implications for Compressor Design}

The three principles above suggest concrete modifications to Memento-style
training.

\textbf{Training signal augmentation.}
Add an auxiliary loss that penalizes mementos that fail to predict,
from the compressed representation alone, what alternative frames were
considered and why they were rejected. This does not require human
annotation: the uncompressed chain-of-thought contains this information,
and the compressor can be trained to preserve it via self-supervised
objectives.

\textbf{SSC-perturbed evaluation.}
Include in the evaluation suite a set of follow-up reasoning contexts
deliberately constructed to resemble the acquisition contexts of
alternative frames that appeared in the compressed blocks. Mementos
that fail on these perturbed contexts should receive lower quality scores
regardless of their in-distribution accuracy.

\textbf{Stack-aware compression architectures.}
Rather than training a single compressor for all reasoning steps,
train separate compressors for steps at different stack depths
(early in reasoning, when multiple frames are still live, vs.\ late,
when commitment has occurred). Early-step compressors should be trained
with stronger alternative-frame index preservation objectives.

\section{Measuring What Survives: An RSA-Based Pilot}
\label{sec:pilot}

Principles 2 and 3 require measuring whether alternative-frame
information survives compaction.
We demonstrate that RSA provides a tractable proxy for this
measurement.

\paragraph{Setup.}
If compaction is a frame-selection operation, then the relational
geometry of the compressed state should reflect which frames survived.
We operationalize this by computing pairwise dissimilarity matrices
$D_{ij} = 1 - \mathrm{corr}(\mathbf{h}_i, \mathbf{h}_j)$ over
$n = 29$ probe concepts at hidden layers of Pythia-70m
\citep{biderman2023} across three training stages: random
initialization (step~0), early training (step~1000), and fully
trained reference (step~143000).
The second-order Spearman correlation $\rho_\mathrm{cog}$ between
checkpoint pairs measures how much structural information is retained
as training progresses.

\paragraph{Results.}
$\rho_\mathrm{cog}$ increases monotonically with training
(L6: step-0 $= 0.017$; step-1000 $= 0.535$), indicating that
training progressively accumulates relational structure.
Crucially, step-0 is not structurally uncorrelated with the
trained reference---early structure persists rather than being
deleted, consistent with the interpretive stack model.
Cross-model comparison (GPT-2 vs.\ Pythia-trained:
$\rho_\mathrm{cog} \approx 0.42$; vs.\ Pythia-step-0:
$\rho_\mathrm{cog} \approx 0.10$) demonstrates that the
protocol discriminates structurally distinct representations.

\paragraph{Relevance to Principle 3.}
Applied to context compression, the same protocol provides
an operationalization of SSC-aware evaluation (Principle 3).
Given a compressed memento $m_t$, one can compute the RSA
dissimilarity matrix of the model conditioned on $m_t$
and compare it to the matrix conditioned on the full
uncompressed context.
A high $\rho_\mathrm{cog}$ indicates that the compression
preserved the relational geometry of the active frames;
a low $\rho_\mathrm{cog}$---especially for concepts
associated with the suppressed alternative frame---indicates
that alternative-frame index information has been lost.
This provides a quantitative signal for the auxiliary loss
proposed in Section~5 without requiring human annotation.
All code is available at
\url{https://github.com/minamominamoto/grounding-without-reality}.

\section{Conclusion}

Context compression is a necessary engineering capability for reasoning
models operating under resource constraints. Memento and related approaches
demonstrate that significant compression ratios are achievable with
acceptable accuracy loss on in-distribution benchmarks. We have argued
that this success is accompanied by a structural risk: compression
optimized on i.i.d.\ accuracy systematically discards the alternative-frame
indices that enable robust reasoning under distributional shift.

The resolution requires not a change to the compression architecture but
a change to the evaluation criterion and training objective: compressors
should be trained and evaluated not only on conclusion-preservation accuracy
but on SSC---the stability of frame selection under the perturbations that
characterize out-of-distribution reasoning contexts.

The three principles proposed here (conclusion preservation, alternative-frame
index preservation, and SSC-aware evaluation) are not technically demanding
in isolation. Their implementation requires primarily a richer theory of
what counts as important to preserve---a question that current engineering
metrics are not equipped to answer.

\bibliographystyle{plainnat}
\begin{thebibliography}{10}

\bibitem[Kontonis et al.(2026)]{kontonis2026memento}
Kontonis, V., et al. (2026).
Memento: Teaching LLMs to manage their own context for long-horizon
reasoning.
arXiv preprint.

\bibitem[LightThinker++(2026)]{lightthinkerplusplus2026}
LightThinker++ Team (2026).
LightThinker++: Commit, Expand, Fold for efficient reasoning.
arXiv preprint.

\bibitem[Neural Computer(2026)]{neuralcomputer2026}
Neural Computer Team (2026).
Neural Computer: Collapsing computation, memory, and I/O.
arXiv preprint.

\bibitem[Lanchantin et al.(2026)]{lanchantin2026midtraining}
Lanchantin, J., et al. (2026).
Thinking LLMs: General instruction following with thought generation
via mid-training.
arXiv:2601.21343.

\bibitem[Minamoto(2026a)]{minamoto2026overwrite}
Minamoto, M. (2026a).
Abduction as overwrite: Negative abductive capability as a structural
condition for reliable inference.
Preprint.
\url{https://github.com/minamominamoto/topology-of-grounding}

\bibitem[Minamoto(2026b)]{minamoto2026accum}
Minamoto, M. (2026b).
Abductive accumulation: Why interpretive frames persist and how
selection stability determines cognitive reliability.
Preprint.
\url{https://github.com/minamominamoto/topology-of-grounding}

\bibitem[Biderman et al.(2023)]{biderman2023}
Biderman, S., et al.\ 2023.
Pythia: A suite for analyzing large language models across training and scaling.
\textit{Proceedings of ICML 2023}.

\end{thebibliography}

\end{document}
